\chapter{Surgery on a $\delta$-neck}
\label{chap:surgery-delta-neck}

\section{Notation and surgery metric}
\label{necksurgery}

Fix a $\delta$-neck $(N,g)$ centered at $x_0$, represented by
\[
  \rho:S^2\times(-\delta^{-1},\delta^{-1})\longrightarrow N,
\]
and the standard initial metric $(\Ar^3,g_0)$. Write $h_0+ds^2$ for the
product of the round scalar-curvature-one metric on $S^2$ and the Euclidean
metric. Let $N^-=\rho(S^2\times(-\delta^{-1},0])$ and use $s$ for its second
coordinate.

Let $p_0$ be the tip of the rotationally symmetric standard metric. Fix
$A_0>0$ and an isometry
\[
  \psi:(S^2\times(-\infty,4],h_0+ds^2)
  \longrightarrow(\Ar^3\setminus B(p_0,A_0),g_0).
\]
The induced coordinate $s_1$ extends continuously to $\Ar^3$ by
$s_1(x)=A_0+4-d_{g_0}(p_0,x)$, is smooth away from $p_0$, and has round
$2$-sphere levels of scalar curvature at least $1$. Choose
$0<r_0<A_0$ so that $g_0$ has constant sectional curvature $1/4$ on
$B(p_0,r_0)$.

Form ${\mathcal S}$ by gluing
$\rho(S^2\times(-\delta^{-1},4))$ to $B(p_0,A_0+4)$ using $\rho(x,s)=\psi(x,s)$
for $0<s<4$. The two radial coordinates define
$s:{\mathcal S}\to(-\delta^{-1},4+A_0]$.

\begin{claim}[Distance-decreasing map from delta-neck to surgery cap model]
\label{claim:neck-to-cap-distance-decreasing}
\label{distdecrease}
Set $\eta=\sqrt{1-\delta}$. Define $\xi:N\to\Ar^3$ by sending
$\rho(S^2\times[A_0+4,\delta^{-1}))$ to $p_0$ and sending
$\rho(x,s)$, for $s<A_0+4$, to the point on the radial ray $x$ at
$g_0$-distance $A_0+4-s$ from $p_0$. Then
\[
  \xi:(N,R(x_0)g)\longrightarrow(\Ar^3,\eta g_0)
\]
is distance nonincreasing.
\end{claim}

\begin{proof}
\sketch
\uses{lem:standard-initial-metric-distance-decreasing}
The neck condition gives
$R(x_0)g\geq\eta(h_0+ds^2)$ for small $\delta$. The standard radial map is
distance nonincreasing from $h_0+ds^2$ to $g_0$ by Lemma~\ref{decrease}; the
claim follows after scaling.
\end{proof}

Fix cutoff functions $\alpha:[1,2]\to[0,1]$ and
$\beta:[4+A_0-r_0,4+A_0]\to[0,1]$, equal to $1$ near their left endpoints
and $0$ near their right endpoints. For constants $C_0,q>0$, define
\[
  f(s)=
  \begin{cases}
  0,&s\leq0,\\
  C_0\delta e^{-q/s},&s>0.
  \end{cases}
\]
Put $A_{r_0}=4+A_0-r_0$ and $A'=4+A_0$, and define
\[
\widehat g=
\begin{cases}
e^{-2f(s)}R(x_0)\rho^*g,&s\leq1,\\
e^{-2f(s)}\bigl(\alpha(s)R(x_0)\rho^*g+(1-\alpha(s))\eta g_0\bigr),
  &1\leq s\leq2,\\
e^{-2f(s)}\eta g_0,&2\leq s\leq A_{r_0},\\
\bigl(\beta(s)e^{-2f(s)}+(1-\beta(s))e^{-2f(A')}\bigr)\eta g_0,
  &A_{r_0}\leq s\leq A'.
\end{cases}
\qquad
\widetilde g=R(x_0)^{-1}\widehat g.
\]

\begin{theorem}[Local surgery on a delta-neck: curvature control after surgery]
\label{thm:local-surgery}
\uses{claim:neck-to-cap-distance-decreasing}

\label{LOCALSURGERY}
There are $C_0,q,R_0<\infty$ and $\delta'_0>0$ such that if
$R(x_0)\geq R_0$ and $0<\delta\leq\delta'_0$, the surgery result
$({\mathcal S},\widetilde g)$ has the following properties.
\begin{enumerate}
\item If $t\geq0$, $X(p)=\max(0,-\nu_g(p))$, and every $p\in N$ satisfies
\[
  R(p)\geq-\frac6{1+4t},
  \qquad
  R(p)\geq2X(p)\bigl(\log X(p)+\log(1+t)-3\bigr)
\]
whenever $X(p)>0$, then the same inequalities hold on
$({\mathcal S},\widetilde g)$ with the same $t$.
\item The metric $\widetilde g$ has positive sectional curvature on
$s^{-1}([1,4+A_0])$.
\item The map $\xi$ of Claim~\ref{distdecrease}, regarded as a map into the
surgery cap, is distance nonincreasing from $g$ to $\widetilde g$.
\item For every $\delta''>0$ there is
$\delta'_1(\delta'')>0$ such that, if
$\delta\leq\min(\delta'_0,\delta'_1)$, then
$\widehat g$ on $B_{\widehat g}(p_0,(\delta'')^{-1})$ is
$\delta''$-close in $C^{[1/\delta'']}$ to $g_0$ on
$B_{g_0}(p_0,(\delta'')^{-1})$.
\end{enumerate}
\end{theorem}

\begin{proof}
\sketch
On $s<4$, Corollaries~\ref{Rcompar} and~\ref{smeigen} show that scalar
curvature does not decrease and the least curvature eigenvalue increases.
Proposition~\ref{prop:pinching-preserved-s-less-4} therefore preserves the two
pinching inequalities, while Lemma~\ref{nuprime>0} gives positivity on
$1\leq s<4$. On $4\leq s\leq4+A_0-r_0$, positivity follows from the
nonnegative curvature of $g_0$ and the same conformal comparison; on the
remaining compact tip region, $\widetilde g\to g_0$ smoothly and $g_0$ has
positive curvature.

For the distance assertion, the old and new metrics agree on $s\leq0$;
on $0\leq s\leq4$ the defining interpolation satisfies
$R(x_0)\rho^*g\geq\widehat g$, and on the cap Claim~\ref{distdecrease}
applies. Finally, $f\to0$ smoothly and $\eta\to1$ as $\delta\to0$, giving the
stated convergence to the standard metric.
\end{proof}

\begin{definition}[Surgery cap]
\label{def:surgery-cap}
\label{surgerycap}
The image of $B_{g_0}(p_0,A_0+4)$ in ${\mathcal S}$ is the
\emph{surgery cap}.
\end{definition}

\begin{lemma}[Surgery cap lies between two concentric metric balls]
\label{lem:surgery-cap-ball-bounds}
\uses{def:surgery-cap}

\label{capball}
For sufficiently small $\delta$ and $\epsilon<1/200$, the surgery cap lies in
\[
  B_{\widetilde g}
  \bigl(p_0,R(x_0)^{-1/2}(A_0+5)\bigr)
\]
and contains
\[
  B_{\widetilde g}
  \bigl(p_0,R(x_0)^{-1/2}(A_0+3)\bigr).
\]
The complement of its closure is isometric to $N^-$.
\end{lemma}

\section{Conformal curvature estimates}

Let $I\subset(-\delta^{-1},4+A_0)$ be open, let $h$ on $S^2\times I$ be
$\delta$-close in $C^{[1/\delta]}$ to $h_0+ds^2$, and set
$\widehat h=e^{-2f}h$. The conformal curvature formula is
\begin{eqnarray}
\label{confchange}
\widehat R_{ijkl}
&=&e^{-2f}\bigl(R_{ijkl}-f_jf_kh_{il}+f_jf_lh_{ik}
+f_if_kh_{jl}-f_if_lh_{jk}\\
&&\quad-(\wedge^2h)_{ijkl}|\nabla f|^2-f_{jk}h_{il}+f_{ik}h_{jl}
+f_{jl}h_{ik}-f_{il}h_{jk}\bigr),\nonumber
\end{eqnarray}
where $f_{ij}=\partial_if_j-f_l\Gamma^l_{ij}$ and
$(\wedge^2h)_{ijkl}=h_{ik}h_{jl}-h_{il}h_{jk}$ \cite{Sakai}.

Use coordinates $(x^0,x^1,x^2)$ with $x^0=s$. The neck estimates give
\[
\begin{aligned}
|\nabla f|_h^2&=\frac{q^2}{s^4}f^2(1+O(\delta)),\\
f_{00}&=\left(\frac{q^2}{s^4}-\frac{2q}{s^3}\right)f
+\frac q{s^2}fO(\delta),\\
f_{i0}&=\frac q{s^2}fO(\delta),
&f_{ij}&=\frac q{s^2}fO(\delta)\quad(1\leq i,j\leq2).
\end{aligned}
\]
Choose $q$, independently of $C_0$ and $\delta$, so that
\begin{equation}
\label{qeqn}
  q\gg(4+A_0)^2,
  \qquad
  \frac{q^2}{s^4}e^{-q/s}\ll1
  \quad(0<s\leq4+A_0).
\end{equation}
After choosing $C_0$, require $\delta\ll C_0^{-1}$; then
\[
  \frac q{s^2}f^2\ll\frac{q^2}{s^4}f^2
  \ll\frac q{s^2}f\ll1.
\]

\begin{corollary}[Curvature formula for the conformally rescaled neck metric]
\label{cor:conformal-curvature-formula}
\label{delta0}
For fixed $C_0,q$ there is $\delta'_2>0$ such that, if
$\delta\leq\delta'_2$, then
\begin{equation}
\label{maineqn}
(\widehat R_{ijkl})=e^{-2f}\left[
(R_{ijkl})+
\begin{pmatrix}
-\dfrac{q^2}{s^4}f^2&0\\
0&\left(\dfrac{q^2}{s^4}-\dfrac{2q}{s^3}\right)fI_2
\end{pmatrix}
+\frac q{s^2}fO(\delta)\right].
\end{equation}
Moreover,
\[
  \widehat R=e^{2f}\left[
  R+4\left(\frac{q^2}{s^4}-\frac{2q}{s^3}\right)f
  -2\frac{q^2}{s^4}f^2+\frac q{s^2}fO(\delta)\right].
\]
\end{corollary}

\begin{proof}
\sketch
Insert the displayed first- and second-derivative estimates for $f$ into
Equation~\eqref{confchange}. The $00$ Hessian term supplies the positive
$2\times2$ block, the gradient-square term supplies the negative scalar block,
and all mixed terms are absorbed into $qfO(\delta)/s^2$. Tracing the conformal
formula gives the scalar-curvature identity.
\end{proof}

\begin{corollary}[Scalar curvature does not decrease under the surgery conformal factor]
\label{cor:scalar-curvature-nondecreasing-conformal}
\uses{cor:conformal-curvature-formula}

\label{Rcompar}
For fixed $C_0<\infty$ and
$\delta<\min(\delta'_2,C_0^{-1})$, one has $\widehat R\geq R$.
\end{corollary}

\begin{proof}
\sketch
Since $C_0\delta<1$, $f^2\ll f$; Equation~\eqref{qeqn} makes the positive
$q^2f/s^4$ term dominate the terms involving $qf/s^3$, $qf/s^2$, and
$q^2f^2/s^4$.
\end{proof}

\begin{lemma}[Near-diagonal curvature eigenbasis for a near-cylindrical metric]
\label{lem:near-cylinder-curvature-eigenbasis}
\label{newbasis}
There is $\delta'_3>0$ such that, if $\delta\leq\delta'_3$ and
$\{e_0,e_1,e_2\}$ is product-orthonormal with $e_0$ radial, then there is an
$h$-orthonormal basis $\{f_0,f_1,f_2\}$ whose change-of-basis matrix is
$I+O(\delta)$ and for which
\[
  (R_{ijkl})=
  \begin{pmatrix}1/2&0&0\\0&0&0\\0&0&0\end{pmatrix}+O(\delta)
\]
in the induced basis of $\wedge^2T_y(S^2\times I)$.
\end{lemma}

\begin{proof}
\sketch
The matrices of $h$ and its curvature in the product frame are respectively
$I+O(\delta)$ and
$\operatorname{diag}(1/2,0,0)+O(\delta)$. Gram--Schmidt changes the frame by
$I+O(\delta)$, and conjugation by its exterior square preserves these error
bounds.
\end{proof}

\begin{corollary}[Curvature block-diagonalizes in the adapted orthonormal basis]
\label{cor:curvature-diagonal-block-basis}
\uses{lem:near-cylinder-curvature-eigenbasis}

\label{2ndnewbasis}
The basis in Lemma~\ref{newbasis} may be chosen so that
\[
  (R_{ijkl})=
  \begin{pmatrix}\lambda&0&0\\0&\alpha&\beta\\0&\beta&\gamma\end{pmatrix},
\]
where
$|\lambda-1/2|,|\alpha|,|\beta|,|\gamma|=O(\delta)$.
\end{corollary}

\begin{proof}
\sketch
The maximum eigendirection of the curvature quadratic form is
$O(\delta)$-close to the first basis vector. Taking it as the first vector
eliminates the mixed first row and column, while the restriction to its
orthogonal complement is $O(\delta)$.
\end{proof}

\begin{lemma}[Curvature matrix of the conformally rescaled metric in the adapted basis]
\label{lem:rhat-matrix-adapted-basis}
\uses{cor:curvature-diagonal-block-basis}

\label{Rhatmatrix}
Let $\delta'_4=\min(\delta'_2,\delta'_3)$. If
$\delta\leq\min(\delta'_4,C_0^{-1})$, then in the basis of
Corollary~\ref{2ndnewbasis},
\[
(\widehat R_{ijkl})=e^{-2f}\left[
\begin{pmatrix}\lambda&0&0\\0&\alpha&\beta\\0&\beta&\gamma\end{pmatrix}
+\begin{pmatrix}
-\dfrac{q^2}{s^4}f^2&0\\
0&\left(\dfrac{q^2}{s^4}-\dfrac{2q}{s^3}\right)fI_2
\end{pmatrix}
+\frac{q^2}{s^4}fO(\delta)\right].
\]
\end{lemma}

\begin{proof}
\sketch
\uses{cor:conformal-curvature-formula}
Conjugate Equation~\eqref{maineqn} by the $I+O(\delta)$ change of basis.
The choice of $q$ and $C_0\delta<1$ absorb the resulting errors into the final
term.
\end{proof}

\begin{corollary}[Fully diagonalized curvature matrices for the neck metric and its rescaling]
\label{cor:diagonalized-curvature-matrices}
\uses{lem:near-cylinder-curvature-eigenbasis, lem:rhat-matrix-adapted-basis}
\label{13.10}

If $\delta\leq\min(\delta'_4,C_0^{-1})$, there is an $h$-orthonormal basis in
which
\[
  (R_{ijkl})=\operatorname{diag}(\lambda,\mu,\nu),
  \qquad
  |\lambda-1/2|,|\mu|,|\nu|=O(\delta),
\]
and
\[
  (\widehat R_{ijkl})=e^{-2f}\left[
  \operatorname{diag}(\lambda,\mu,\nu)
  +\operatorname{diag}\left(-\frac{q^2}{s^4}f^2,
  \left(\frac{q^2}{s^4}-\frac{2q}{s^3}\right)f,
  \left(\frac{q^2}{s^4}-\frac{2q}{s^3}\right)f\right)
  +\frac{q^2}{s^4}fO(\delta)\right].
\]
\end{corollary}

\begin{proof}
\sketch
\uses{lem:near-cylinder-curvature-eigenbasis}
An orthogonal rotation in the last two exterior-square directions diagonalizes
the lower block without changing the asymptotic bounds or the scalar form of
the conformal correction on that block.
\end{proof}

\begin{corollary}[Eigenvalue estimates for the conformally rescaled curvature]
\label{cor:conformal-eigenvalue-estimates}
\label{improve}
For the chosen $q$ there is $A<\infty$ such that, for every fixed $C_0$ and
sufficiently small $\delta$, if
$\lambda\geq\mu\geq\nu$ are the curvature eigenvalues of $h$ and
$\lambda',\mu',\nu'$ those of $\widehat h$, then
\[
\begin{aligned}
\left|\lambda'-e^{2f}\left(\lambda-\frac{q^2}{s^4}f^2\right)\right|
&\leq\frac{q^2}{s^4}fA\delta,\\
\left|\mu'-e^{2f}\left(\mu+left(\frac{q^2}{s^4}-\frac{2q}{s^3}\right)f\right)\right|
&\leq\frac{q^2}{s^4}fA\delta,\\
\left|\nu'-e^{2f}\left(\nu+left(\frac{q^2}{s^4}-\frac{2q}{s^3}\right)f\right)\right|
&\leq\frac{q^2}{s^4}fA\delta.
\end{aligned}
\]
In particular,
\[
  \mu'\geq e^{2f}\left(\mu+\frac{q^2}{2s^4}f\right),
  \qquad
  \nu'\geq e^{2f}\left(\nu+\frac{q^2}{2s^4}f\right).
\]
\end{corollary}

\begin{proof}
\sketch
\uses{cor:diagonalized-curvature-matrices,lem:rhat-matrix-adapted-basis}
The frame $\{e^ff_i\}$ is orthonormal for $\widehat h$, so the matrix in
Corollary~\ref{13.10} gains the factor $e^{4f}$ and hence has leading factor
$e^{2f}$. Standard symmetric-matrix perturbation bounds give the first three
estimates. The choice $q\gg(4+A_0)^2$ and small $\delta$ yield the final two.
\end{proof}

\begin{corollary}[Smallest curvature eigenvalue increases under the surgery conformal change]
\label{cor:smallest-eigenvalue-increases}
\uses{cor:conformal-eigenvalue-estimates}

\label{smeigen}
For the chosen $q$, every fixed $C_0$, and sufficiently small $\delta$, the
least eigenvalue of $\Rm_{\widehat h}$ is greater than that of $\Rm_h$ at the
same point. Thus nonnegative curvature of $h$ implies nonnegative curvature of
$\widehat h$.
\end{corollary}

\begin{proof}
\sketch
Because $\lambda=1/2+O(\delta)$ and $\mu,\nu=O(\delta)$, while
$q^2f/s^4\ll1$, the least new eigenvalue is $\mu'$ or $\nu'$. The final
inequalities of Corollary~\ref{improve} make each strictly larger than its old
counterpart.
\end{proof}

Let $K$ be a universal constant such that every curvature eigenvalue of a
scale-one $\delta$-neck is at least $-K\delta$, and set
$C_0=2Ke^q$.

\begin{lemma}[Positivity of curvature eigenvalues on the surgery region]
\label{lem:positive-eigenvalues-surgery-region}
\uses{cor:conformal-eigenvalue-estimates}

\label{nuprime>0}
For the chosen $q,C_0$ and sufficiently small $\delta$, on
$1\leq s\leq4+A_0$ one has
\[
  \mu'>0,\qquad\nu'>0,\qquad\lambda'>\frac14.
\]
\end{lemma}

\begin{proof}
\sketch
The function $q^2f/(2s^4)$ is increasing on $[1,4+A_0]$ and at $s=1$ equals
$(q^2/2)e^{-q}C_0\delta>K\delta$. It therefore dominates the lower bound
$-K\delta$ for both small eigenvalues. The correction
$q^2f^2/s^4$ is arbitrarily small, so $\lambda'=1/2+o(1)>1/4$.
\end{proof}

\begin{proposition}[Curvature pinching is preserved by surgery on the region $s<4$]
\label{prop:pinching-preserved-s-less-4}
For sufficiently small $\delta$, suppose the curvature of the original neck
satisfies, for some $t\geq0$,
\[
  R\geq-\frac6{1+4t},
  \qquad
  R\geq2X\bigl(\log X+\log(1+t)-3\bigr)
  \quad(X>0).
\]
Then $({\mathcal S},\widetilde g)$ satisfies the same inequalities on
$s^{-1}(( -\delta^{-1},4))$ and has positive curvature on
$s^{-1}([1,4))$.
\end{proposition}

\begin{proof}
\sketch
\uses{thm:local-surgery,cor:scalar-curvature-nondecreasing-conformal,lem:positive-eigenvalues-surgery-region,cor:smallest-eigenvalue-increases}
Corollaries~\ref{smeigen} and~\ref{Rcompar} give
$X_{\widehat h}\leq X_h$ and $\widehat R\geq R$. If
$X_h\geq e^3/(1+t)$, the function
$u\mapsto2u(\log u+\log(1+t)-3)$ is increasing, so the pinching inequality
passes to $\widehat h$. If $X_h<e^3/(1+t)$, neck closeness and small $\delta$
give $\widehat R\geq R\geq0$, whereas the required right-hand side is
negative. Lemma~\ref{nuprime>0} supplies positivity on $1\leq s<4$.
\end{proof}

\section{Other metric control}

\begin{lemma}[Metric balls in the surgery result compared to the neck coordinate]
\label{lem:surgery-result-metric-ball-bounds}
\label{stdballs}
For sufficiently small $\delta$, let $({\mathcal S},\widetilde g)$ be obtained
by surgery on a $\delta$-neck. For every $D>0$, the boundary of
$B_{\widetilde g}(p_0,D+5+A_0)$ is contained in
\[
  s_N^{-1}(-(2D+2),-D/2).
\]
\end{lemma}

\begin{proof}
\sketch
The surgery manifold is naturally identified with the standard ball
$B_{g_0}(p_0,A_0+4+\delta^{-1})$, and the pulled-back surgery metric is
$2\delta$-close to $g_0$. Radial distance comparison in this identification
places the indicated sphere in the stated neck-coordinate annulus.
\end{proof}
